\section{Introduction}
\label{sec:introduction}

% P1: background -> capability gap -> importance.
[The field aims to achieve a long-term goal]. Current large language models (LLMs) can [established capability], while \CoreTerm{} requires [bounded additional capability] in [specific setting]. This distinction matters because [task-grounded consequence]. The overview in \cref{fig:teaser} connects this research question to the proposed design and its eventual evidence.

% P2: classify prior work; put citations beside the supported objects.
Existing work can be organized by [task resources], [model design], and [evaluation]. For example, the Transformer~\citep{vaswani2017attention} is a verified citation example for an architecture name. Replace this syntax demonstration with sources relevant to the actual paper and explain which objective each research direction serves.

% P3: parallel challenges, then an integrating sentence.
Three challenges define the study. First, \textbf{[Challenge A]} concerns [capability and source of difficulty]. Second, \textbf{[Challenge B]} concerns [capability and failure mechanism]. Third, \textbf{[Challenge C]} concerns [capability and evaluation consequence]. Together, \CoreTerm{} requires a design that connects these capabilities through a common task and evaluation protocol.

% P4: same order as the challenge paragraph.
In this paper, we introduce [approach] for \CoreTerm{} under [scope]. [Component A] addresses [Challenge A] by [mechanism]. [Component B] addresses [Challenge B] through [constraint]. [Component C] addresses [Challenge C] with [verification or feedback]. The formulation, construction, and evaluation are developed in \cref{sec:method}.

% P5: verified observations first, interpretation second; no invented numbers.
Experiments on [setting] show [verified quantitative observation] and [verified qualitative observation]. These findings suggest [bounded capability insight], supported by the main results, ablations, and diagnostic cases in \cref{sec:experiments}. Populate this paragraph only after the evidence is available.

% P6: compactitem; do not put an unescaped [label] immediately after \item.
Our contributions are organized around the research question:
\begin{compactitem}
\item \textbf{Task formulation.} [Define the capability, inputs, outputs, and scope.]
\item \textbf{Approach or resource.} [State the method or dataset contribution at a research level.]
\item \textbf{Evaluation and evidence.} [State the evaluation design and supported finding.]
\end{compactitem}
